\uuid{2qL9}
\titre{Différentiabilité en un point }

\niveau{}
\module{ Fonctions de plusieurs variables }
\chapitre{Continuité et différentiabilité}
\sousChapitre{}
\theme{}
\auteur{Maxime Nguyen}
\datecreate{ 2026-04-07}
\organisation{AMSCC}
\difficulte{}

\contenu{

\texte{ 
On considère la fonction $f:\mathbb{R}^2\to\mathbb{R}$ définie par
\[
f(x,y)=
\begin{cases}
\dfrac{x^3+y^3}{x^2+y^2} & \text{si } (x,y)\neq(0,0),\\[0.5em]
0 & \text{si } (x,y)=(0,0).
\end{cases}
\]
}

\begin{enumerate}

\item   \question{Montrer que $f$ est continue en $(0,0)$.}
\indication{}
\reponse{
Posons
\[
r=\sqrt{x^2+y^2}.
\]
Pour $(x,y)\neq(0,0)$, on a $|x|\le r$ et $|y|\le r$, donc
\[
|x^3+y^3|\le |x|^3+|y|^3\le r^3+r^3=2r^3.
\]
Comme $x^2+y^2=r^2$, on obtient
\[
|f(x,y)|=\left|\frac{x^3+y^3}{x^2+y^2}\right|
\le \frac{2r^3}{r^2}=2r.
\]
Lorsque $(x,y)\to(0,0)$, on a $r\to 0$, donc
\[
|f(x,y)|\to 0.
\]
Ainsi
\[
f(x,y)\to 0=f(0,0),
\]
ce qui montre que $f$ est continue en $(0,0)$.
}

\item   \question{Étudier l’existence des dérivées partielles
\[
\frac{\partial f}{\partial x}(0,0)
\qquad \text{et} \qquad
\frac{\partial f}{\partial y}(0,0),
\]
et les calculer si elles existent.}
\indication{}
\reponse{
Par définition,
\[
\frac{\partial f}{\partial x}(0,0)
=
\lim_{h\to 0}\frac{f(h,0)-f(0,0)}{h}.
\]
Or, pour $h\neq 0$,
\[
f(h,0)=\frac{h^3}{h^2}=h.
\]
Donc
\[
\frac{f(h,0)-f(0,0)}{h}=\frac{h-0}{h}=1,
\]
et par suite
\[
\frac{\partial f}{\partial x}(0,0)=1.
\]

De même,
\[
\frac{\partial f}{\partial y}(0,0)
=
\lim_{h\to 0}\frac{f(0,h)-f(0,0)}{h}.
\]
Or, pour $h\neq 0$,
\[
f(0,h)=\frac{h^3}{h^2}=h.
\]
Donc
\[
\frac{f(0,h)-f(0,0)}{h}=\frac{h-0}{h}=1,
\]
d’où
\[
\frac{\partial f}{\partial y}(0,0)=1.
\]

Les deux dérivées partielles existent donc, et
\[
\nabla f(0,0)=(1,1).
\]
}

%\item   \question{Déterminer le gradient de $f$ en tout point de $\mathbb{R}^2\setminus\{(0,0)\}$.}
%\indication{}
%\reponse{
%Pour $(x,y)\neq(0,0)$, on pose
%\[
%N(x,y)=x^3+y^3,\qquad D(x,y)=x^2+y^2.
%\]
%Alors
%\[
%f(x,y)=\frac{N(x,y)}{D(x,y)}.
%\]

%Par la formule de dérivation d’un quotient,
%\[
%\frac{\partial f}{\partial x}(x,y)
%=
%\frac{3x^2(x^2+y^2)-(x^3+y^3)\,2x}{(x^2+y^2)^2}.
%\]
%En développant,
%\[
%\frac{\partial f}{\partial x}(x,y)
%=
%\frac{3x^4+3x^2y^2-2x^4-2xy^3}{(x^2+y^2)^2}
%=
%\frac{x^4+3x^2y^2-2xy^3}{(x^2+y^2)^2}.
%\]

%De même,
%\[
%\frac{\partial f}{\partial y}(x,y)
%=
%\frac{3y^2(x^2+y^2)-(x^3+y^3)\,2y}{(x^2+y^2)^2},
%\]
%d’où
%\[
%\frac{\partial f}{\partial y}(x,y)
%=
%\frac{3x^2y^2+3y^4-2x^3y-2y^4}{(x^2+y^2)^2}
%=
%\frac{-2x^3y+3x^2y^2+y^4}{(x^2+y^2)^2}.
%\]

%Ainsi, pour $(x,y)\neq(0,0)$,
%\[
%\nabla f(x,y)
%=
%\left(
%\frac{x^4+3x^2y^2-2xy^3}{(x^2+y^2)^2},
%\frac{-2x^3y+3x^2y^2+y^4}{(x^2+y^2)^2}
%\right).
%\]
%}

\item   \question{Montrer que $f$ n'est pas différentiable en $(0,0)$. La fonction $f$ est-elle de classe $C^1$ en $(0,0)$ ?}
\indication{}
\reponse{
Si $f$ était différentiable en $(0,0)$, sa différentielle serait donnée par
\[
df_{(0,0)}(h,k)=\frac{\partial f}{\partial x}(0,0)h+\frac{\partial f}{\partial y}(0,0)k=h+k,
\]
puisque les dérivées partielles en $(0,0)$ valent toutes deux $1$.

Il faudrait donc avoir
\[
f(x,y)-(x+y)=o\!\left(\sqrt{x^2+y^2}\right).
\]

Étudions cette quantité le long de la droite $y=x$, avec $x\neq 0$. On a
\[
f(x,x)=\frac{x^3+x^3}{x^2+x^2}=\frac{2x^3}{2x^2}=x.
\]
Donc
\[
f(x,x)-(x+x)=x-2x=-x.
\]
Par conséquent,
\[
\frac{|f(x,x)-2x|}{\sqrt{x^2+x^2}}
=
\frac{|x|}{\sqrt{2}\,|x|}
=
\frac{1}{\sqrt2}.
\]
Cette quantité ne tend pas vers $0$ lorsque $x\to 0$. Donc $f$ n’est pas différentiable en $(0,0)$.

En particulier, une fonction de classe $C^1$ est différentiable. Comme $f$ n’est pas différentiable en $(0,0)$, elle n’est pas de classe $C^1$ en $(0,0)$.
}

\end{enumerate}

}